\documentclass{homework}
% \usepackage{graphicx} % Required for inserting images
\usepackage{amsmath}
\usepackage{gensymb}
\usepackage{subfig}
\usepackage{float}

\class{ECEN 5738}
\title{Homework 4}
\author{Sean Jennings}
\date{March 3, 2024}

\begin{document}
\maketitle

\section{Problem 1}
\begin{letterenum}
\item Taking the derivative w.r.t. time of $W$, we have:
\begin{equation*}
    \frac{dW}{dt} = x^T \dot{x} + \dot{x}^T x = 2 x^T \dot{x} = 2 x^T f(x)
\end{equation*}
Taking the magnitude of $dW/dt$ and using the Cauchy=Schwatz inequality, we have:
\begin{align*}
    \biggl|\frac{dW}{dt}\biggr| = |2 x^T f(x) | &\leq 2 \normtwo{x} \normtwo{f(x)} \\
        &\leq 2 \normtwo{x} (L \normtwo{x}) \\
        &= 2L \normtwo{x}^2 \\
        &= 2L W(t)
\end{align*}

\item We start by proving the upper bound on $\normtwo{x(t)}$.
Since $\dot{W}(t) \leq 2LW(t)$ implies that $\dot{W}(t) + (-2L)W(t) \leq 0$,
we can apply the lemma from lecture to find an exponential bound for $W(t)$:
\begin{equation*}
    W(t) \leq W(0) e^{2Lt}
\end{equation*}
Substituting the expression for $W(t) = \normtwo{x(t)}^2$ and taking the square root yields
an upper bound for the 2-norm of $x(t)$:
\begin{equation*}
    \normtwo{x} \leq \normtwo{x(0)} e^{Lt}
\end{equation*}

The lemma can also be used to find a lower bound. Specifically, we will show
 \begin{equation}
    \dot{W}(t) + \alpha W(t) \geq 0
    \label{eq:lower-bound-assumption}
 \end{equation}
 implies
 \begin{equation*}
    W(t) \geq W(0) e^{-\alpha t}.
 \end{equation*}

To see this, put $Y = -W$ and substitute into eq (\ref{eq:lower-bound-assumption}):
\begin{equation*}
    -\dot{Y}(t) + \alpha (-Y(t)) \geq 0,
\end{equation*}
multiply by -1:
\begin{equation*}
    \dot{Y}(t) + \alpha Y(t) \leq 0,
\end{equation*}
apply the original lemma:
\begin{equation*}
    Y(t) \leq Y(0) e^{-\alpha t},
\end{equation*}
and multiply again by -1:
\begin{equation*}
    W(t) = -Y(t) \geq -Y(0) e^{-\alpha t} = W(0) e^{-\alpha t}.
\end{equation*}

We then follow the same steps as for the upper bound, and see that:
$\dot{W}(t) \geq -2LW(t)$ implies
\begin{equation*}
    \normtwo{x(t)} \geq \normtwo{x(0)} e^{-Lt}
\end{equation*}
\end{letterenum}

\section{Problem 2}

\begin{letterenum}
\item Let $D\subset \R^n$ and $V: D \to \R$ be given by $V(x) = 1/2(x_1^2 + x_2^2)$.
Since $V$ is positive definite, to show that the origin is asymptotically stable,
we will show that there exists a domain $D$ over which $\dot{V}$ is negative definite.
Taking the derivative, we have
    \begin{align*}
        \dot{V}(x) &= x_1 \dot{x}_1 + x_2 \dot{x_2} \\
            &= x_1(-x_1 + x_2^2) + x_2( -2x_2 + 3x_1^2) \\
            &= -x_1^2 + x_1 x_2^2 + -2x_2^2 + 3x_2x_1^2 \\
            &= x_1^2(3x_2 - 1) + x_2^2(x_1 - 2)
    \end{align*}

Thus, $\dot{V}(x) < 0$ when $x_2 < 1/3$ and $x_1 < 2$ (and $x$ is nonzero).
Taking the domain of $V$ to be $D = \{[x_1, x_2] \in \R^2: x_1 < 2, x_2 < 1/3\}$
makes $V$ a negative definite function. Furthermore, the domain $D$ includes the
origin. The origin is therefore asymptotically stable.

\item We assume a generic quadratic form of the Lyapunov candidate function
 and find constants which satisfy the assumptions of Lyapunov's stability theorem.
 Let $V: D \to \R^2$ be
\begin{equation*}
    V(x) = 1/2(ax_1^2 + 2bx_1 x_2 + cx_2^2)
\end{equation*}
for some constants $a,b,c\in\R$. Then, the time derivative is given by
\begin{align*}
    \dot{V}(x) &= ax_1 \dot{x}_1 + b(x_1 \dot{x}_2 + x_2 \dot{x}_1) + c x_2 \dot{x}_2 \\
        &= (ax_1 + b x_2) \dot{x}_1 + (cx_2 + bx_1) \dot{x}_2 \\
        &= (ax_1 + bx_2)(-2x_1 -3x_2 + x_1^2) + (cx_2 + bx_1)(x_1 + x_2) \\
        &= -2ax_1^2 - 3ax_1 x_2 + ax_1^3 - 2bx_1x_2 - 3bx_2^2 + bx_2 x_1^2 \\
        &\qquad + cx_1x_2 + cx_2^2 + bx_1^2 + bx_1 x_2 \\
        &= (-2a + b + ax_1 + bx_2) x_1^2 + (-3a - b + c) x_1 x_2 + (-3b + c) x_2^2
\end{align*}

From the final result of this algebraic manipulation, we find sufficient conditions
for $\dot{V}$ to be negative definite: find values of $a,b$ and $c$ which
make coefficient of the $x_1x_2$ cross-term 0, and the coefficients of the
$x_1^2$ and $x_2^2$ terms negative. That is,
\begin{align}
    -3a - b + c = 0 \label{eq:cross-term} \\
    -2a + b + ax_1 + bx_2 < 0 \label{eq:x_1-term} \\
    -3b + c < 0.  \label{eq:x_2-term}
\end{align}

To satisfy inequality (\ref{eq:x_1-term}), we take the $D\subset \R^2$,
the domain of $V$, to be the half space defined by:
\begin{equation*}
    D = \{[x_1, x_2]\in\R^2: ax_1 + bx_2 < 2a - b \}
\end{equation*}
But this domain must include the origin, so we have $2a - b > 0$, or $a > b/2$.

From equation (\ref{eq:cross-term}), we have that $c = 3a + b$. Substituting 
into inequality (\ref{eq:x_2-term}) yields:
\begin{equation*}
    -3b + 3a + b < 0 \Rightarrow a < 2/3 b.
\end{equation*}

In summation, if $a,b$, and $c$ satisfy
\begin{equation*}
    1/2 b < a < 2/3 b \qquad c = 3a + b,
\end{equation*}
then $\dot{V}$ will be negative definite over $D$. We arbitrarily select
\begin{equation*}
    a = 7 \qquad b = 12 \qquad c = 33
\end{equation*}
which satisfy these relations. Finally, to ensure that the origin is asymptotically
stable, we check that $V$ is positive definite. The expression for $V(x)$ can be
rewritten
\begin{equation*}
    V(x) = \frac{1}{2} x^T \mat{
        a & b \\
        b & c
    } x = \frac{1}{2} x^T \mat{
        7 & 12 \\
        12 & 33
    } x
\end{equation*}
We numerically compute the eigenvalues of this matrix to be
$\lambda_1 \approx 2.308$ and $\lambda_2 \approx 37.692$. Since $\lambda_{1,2} > 0$,
$V$ is positive definite and we have shown the origin in asymptotically stable.

\end{letterenum}

\section{Problem 3}
\begin{letterenum}
\item The non-zero equilibrium point $(x^*, y^*)$ occurs when the first terms of
each of the state equation are zero:
\begin{align*}
    a - by^* &= 0 \\
    cx^* - d &= 0
\end{align*}
Therefore,
\begin{align*}
    x^* = d/c &= 2 \\
    y^* = a/b &= 1 \\
\end{align*}

\item Let $\delta x\in\R^2$ the state variable with coordinates shifted to
the non-zero equilbrium point, given by:
\begin{equation*}
    \delta x = \mat{
        x - x^* \\ y - y^*
    }
\end{equation*}
Then, the linearized dynamics are given by:
\begin{equation*}
    \frac{d}{dt} (\delta x) \approx
        \frac{\delta f}{\delta x} \eval_{(x^*, y^*)} \delta x
\end{equation*}

where $\delta f/\delta x: \R^2 \to \R^2$ is Jacobian of the state equations and
is given by:
\begin{equation*}
    \frac{\delta f}{\delta x}(x, y) = \mat{
        a - by & -bx \\
        cy & cx - d
    }.
\end{equation*}
Evaluating at the equilibrium point, we have
\begin{equation*}
    \frac{\delta f}{\delta x}\eval_{(x^*, y^*)} = \mat{
        0 & -db/c \\
        ca/b & 0
    }
\end{equation*}
We solve the characteristic polynomial
\begin{equation*}
    \lambda^2 + (\frac{ca}{b})(\frac{bd}{c}) =
        \lambda^2 + ad = 0
\end{equation*}
to find the eigenvalues of the linearization are $\lambda_{1,2} = \pm i \sqrt{ad}$.
Since both eigenvalues lie on the imaginary axis, Lyapunov's linearization method
can neither conclude that the nonzero equilibrium point is stable nor unstable.

\item Evaluating the Jacobian at the origin we have:
\begin{equation*}
    \frac{\delta f}{\delta x}\eval{(0, 0)} = \mat{
        a & 0 \\
        0 & -d
    }
\end{equation*}
Thus, the eigenvalues of the linearized system are $\lambda_1 = a$ and 
$\lambda_2 = -d$. Since at least one eigenvalue has a positive real component
($\lambda_1 = a = 1 > 0$), the origin is an unstable equilibrium point.
    
\end{letterenum}

\section{Problem 4}

Let $V(x) = x^T P x$. Since $P$ is positive definite, $V$ is a positive
definite function. Taking the derivative, we have
\begin{align*}
    \dot{V}(x) &= \dot{x}^T P x + x^T P x \\
        &= x^TA^T P x + x^T PA x \\
        &= x^T (A^T P + PA) x \\
        &= x^T (-vv^T) x \\
        &= - (x^T v) (v^T x) \\
        &= - (x^T v)^2
\end{align*}
So $\dot{V}(x) \leq 0$ for all $x\in\R^n$
and $\dot{V}(x) = 0$ if and only if $x^T v = 0$.

Therefore, $\dot{V}$ is negative semidefinite and we can use the LaSalle-Kravoskii
Invariance Principle to show that trajectories converge to the largest invariant
subset of the zero set of $\dot{V}$.

We first define $\Omega = \{x \in \R^n: V(x) \leq c\}$ ($c \in\R$) to be some sublevel set 
of $V$. Since $V$ is nonincreasing along trajectories, it is positively
invariant.  Let $\mathcal{S} = \{x\in\Omega: x^T v = 0\}$ be the set of all $x\in\Omega$ where
$V(x) = 0$. 

Suppose $M \subseteq \mathcal{S}$ is the largest invariant subset of $\mathcal{S}$.
By the invariance principle, any trajectory starting
in $\Omega$ must converge to $M$, so it remains to show that $M$ contains only
the origin.

Select any $x_0 \in M$. Since $M$ is invariant, a trajectory $x(t)$
starting at $x_0$ stays in $M$, (i.e. $x(t) \in M$ for all $t$). 
The system is linear, so it has solution 
\begin{equation*}
    x(t) = e^{At} x_0
\end{equation*}
for all $t$.  But $x(t)$ must also be a zero of $V$, so
\begin{equation*}
    0 = v^T x(t) = v^T e^{At} x_0 
\end{equation*}
and the given assumption implies that $x_0 = 0$.
Since $x_0$ was arbitrary, $M$ contains only the origin, and the origin
is asymptotically stable.


\end{document}
